\section{Limitations}

The first limitation is scope. The strongest results are on audited executable
domains and naturalized/model-generated transfer built around those domains.
These results should not be read as broad open-domain generality.

The second limitation is that the benchmark story remains layered rather than
singular. \textsc{Craft-Clean} is usable at mid-scale under the
artifact-focused acceptance rule, but the external human blind audit is still
pending. \textsc{ProcessBench} is now an operational external-source package
that is closed for its intended corroboration role, but it is not a
matched-quartet replacement and does not by itself reproduce the full
\amcd/\ass{} protocol. We therefore avoid language that implies a single fully
closed benchmark, or independence from this layered package structure.

The third limitation is calibration. The downstream evaluation includes
\texttt{ECE / Brier / AURC}, but the strongest improvements are on ranking and
exploitability rather than probability calibration. Calibration should
therefore remain a secondary downstream result and is reported only in the
appendix.

Finally, the current dual-head and replay lines are informative negative
results rather than positive method stories. They help define the
faithfulness-utility boundary, but they do not currently support a claim that a
single repair mechanism already solves answer leakage at this scale. More
generally, the strongest contribution is an audit-and-deployment result, not an
architecture result.
